\chapter{ארכיטקטורת הטרנספורמר}

ארכיטקטורת ה-\textenglish{Transformer} \cite{vaswani2017attention} מהווה את
אחת מפריצות הדרך המשמעותיות ביותר בתחום עיבוד השפה הטבעית בעשור האחרון.
בשונה מארכיטקטורות הרקורנטיות שקדמו לה, כגון \textenglish{LSTM} ו-\textenglish{GRU},
ה-\textenglish{Transformer} מעבד את כל אסימוני הרצף במקביל, ומסתמך
אך ורק על מנגנוני קשב לצורך מידול התלויות בין אסימונים.
פרק זה בוחן את יסודות הארכיטקטורה: מנגנון הקשב הבסיסי, מנגנון הקשב
הרב-ראשי, מבנה הקידוד-פענוח, וקידוד המיקומי — ארבעת אבני הבניין
שעליהן מושתת כלל מגוון מודלי השפה המודרניים.

\section{המנגנון הבסיסי: שאילתות, מפתחות וערכים}

מנגנון הקשב מוסקאלרי (\textenglish{Scaled Dot-Product Attention}) \cite{vaswani2017attention}
מגדיר פונקציה המקבלת שלוש מטריצות כקלט: מטריצת שאילתות
\textenglish{(Queries)} \(\mathbf{Q} \in \mathbb{R}^{n \times d_k}\),
מטריצת מפתחות \textenglish{(Keys)} \(\mathbf{K} \in \mathbb{R}^{m \times d_k}\),
ומטריצת ערכים \textenglish{(Values)} \(\mathbf{V} \in \mathbb{R}^{m \times d_v}\),
כאשר \(n\) הוא מספר השאילתות, \(m\) הוא מספר זוגות המפתח-ערך,
\(d_k\) הוא ממד המפתח, ו-\(d_v\) הוא ממד הערך.

הפונקציה עצמה מוגדרת כדלקמן:

\begin{equation}
  \text{Attention}(\mathbf{Q}, \mathbf{K}, \mathbf{V}) =
  \text{softmax}\!\left(\frac{\mathbf{Q}\mathbf{K}^{\top}}{\sqrt{d_k}}\right)\mathbf{V}
\end{equation}

ניתן לפרש נוסחה זו באופן הבא: עבור כל שאילתה \(\mathbf{q}_i\),
מחושב מכפלת הנקודות שלה עם כל מפתח \(\mathbf{k}_j\), המייצגת את
מידת הרלוונטיות של האסימון ה-\(j\) לשאילתה ה-\(i\).
ציוני הדמיון מחולקים בגורם הנרמול \(\sqrt{d_k}\) כדי למנוע
ריכוז ערכים גדולים שיגרמו לפונקציית ה-\textenglish{softmax}
להגיע לאזורים רוויים בעלי גרדיאנטים זעירים.
לאחר מכן מופעלת פונקציית ה-\textenglish{softmax} לאורך ממד המפתחות,
המניבה התפלגות הסתברותית המשקפת את משקלי הקשב,
ולבסוף מחושבת ממוצע משוקלל של וקטורי הערכים.

ההצדקה האינטואיטיבית לגישה זו נשענת על מושג ה-\textenglish{soft retrieval}:
בניגוד לחיפוש ב-\textenglish{hash table} שמחזיר ערך בודד המתאים למפתח מדויק,
מנגנון הקשב מחשב סכום רציף של כל הערכים, משוקלל לפי מידת ההתאמה
בין השאילתה לכל מפתח.
תכונה זו מאפשרת למודל לשלב מידע ממקורות מרובים בצורה דיפרנציאלית,
ולפיכך כל הפרמטרים מתעדכנים ביעילות במהלך הלמידה.

ראוי לציין כי בניסוחו המקורי של וסמאני ועמיתיו \cite{vaswani2017attention},
מטריצות \(\mathbf{Q}\), \(\mathbf{K}\) ו-\(\mathbf{V}\) אינן נלקחות ישירות
מייצוגי הקלט, אלא מתקבלות מהם באמצעות טרנספורמציות לינאריות
\(\mathbf{W}^Q\), \(\mathbf{W}^K\) ו-\(\mathbf{W}^V\) הנלמדות במהלך האימון.
הפרמטריזציה הזו מקנה למודל גמישות רבה, שכן ניתן ללמוד מרחבי השאילתות,
המפתחות והערכים בנפרד, ובכך לייצג תלויות מורכבות בין אסימונים.


\section{קשב רב-ראשי ומבנה קידוד-פענוח}

מנגנון הקשב הרב-ראשי (\textenglish{Multi-Head Attention}) מרחיב את
קשב הנקודתי הבסיסי על ידי הפעלת מספר ראשי קשב במקביל,
כאשר כל ראש לומד מרחב ייצוג שונה ומוקרן באופן עצמאי.
באופן פורמלי, עבור \(h\) ראשי קשב, כל ראש \(i\) מחשב:

\begin{equation}
  \text{head}_i = \text{Attention}\!\left(\mathbf{Q}\mathbf{W}_i^Q,\,
  \mathbf{K}\mathbf{W}_i^K,\,
  \mathbf{V}\mathbf{W}_i^V\right)
\end{equation}

כאשר \(\mathbf{W}_i^Q \in \mathbb{R}^{d_{\text{model}} \times d_k}\),
\(\mathbf{W}_i^K \in \mathbb{R}^{d_{\text{model}} \times d_k}\),
ו-\(\mathbf{W}_i^V \in \mathbb{R}^{d_{\text{model}} \times d_v}\)
הן מטריצות הטרנספורמציה הנלמדות לראש ה-\(i\).
הפלטים של כל הראשים מחוברים (\textenglish{concatenated}) ומוקרנים
באמצעות מטריצה לינארית נוספת \(\mathbf{W}^O\):

\begin{equation}
  \text{MultiHead}(\mathbf{Q}, \mathbf{K}, \mathbf{V}) =
  \text{Concat}\!\left(\text{head}_1, \ldots, \text{head}_h\right)\mathbf{W}^O
\end{equation}

היתרון של ריבוי הראשים הוא מהותי: ראשים שונים יכולים לתפוס
יחסים שונים בין אסימונים — ראש אחד עשוי להתמקד ביחסים
תחביריים, ואחר ביחסים סמנטיים מרחוקים.
מחקרים מאוחרים \cite{clark2020electra} הראו כי ראשי הקשב של
מודלים כגון \textenglish{BERT} \cite{devlin2019bert} מפתחים
ייצוגים ייחודיים לתפקידים בלשניים שונים.

\subsection*{מבנה הקידוד-פענוח}

מבנה הקידוד-פענוח (\textenglish{Encoder-Decoder}) של
ה-\textenglish{Transformer} המקורי \cite{vaswani2017attention}
מורכב משתי ערמות: מקודד (\textenglish{encoder}) ומפענח
(\textenglish{decoder}), כל אחד מהם מורכב ממספר שכבות זהות.

כל שכבת מקודד מכילה שני תת-רכיבים עיקריים:
תת-שכבת קשב עצמי (\textenglish{self-attention}) המאפשרת
לכל אסימון לתשומת לב לכל אסימון אחר ברצף הקלט,
ותת-שכבת רשת נוירונים צפה (\textenglish{Feed-Forward Network}).
כל אחד מתת-הרכיבים עוטף בחיבור שיורי
(\textenglish{residual connection}) ונרמול שכבה
(\textenglish{layer normalization}):

\begin{equation}
  \mathbf{x}_{l+1} = \text{LayerNorm}\!\left(\mathbf{x}_l +
  \text{Sublayer}(\mathbf{x}_l)\right)
\end{equation}

החיבור השיורי, שנשאל מארכיטקטורת \textenglish{ResNet},
מאפשר לגרדיאנטים לזרום ישירות לשכבות עמוקות יותר במהלך
האחור-הפצה (\textenglish{backpropagation}), וכך מקל על
אימון רשתות עמוקות.
נרמול השכבה מייצב את ההתפלגות של ההפעלות בכל שכבה,
ומואץ את ההתכנסות \cite{vaswani2017attention}.

שכבת המפענח, לעומת זאת, מוסיפה תת-שכבה שלישית: קשב
צולב (\textenglish{cross-attention}), שבו השאילתות
נגזרות מהמפענח אך המפתחות והערכים נגזרים ממוצא המקודד.
מנגנון זה מאפשר למפענח לשים לב לחלקים רלוונטיים
ברצף הקלט בעת יצירת כל אסימון בפלט.
בנוסף, שכבת הקשב העצמי של המפענח היא \textenglish{masked}:
היא חוסמת גישה לאסימונים עתידיים על מנת לשמר את
התכונה האוטורגרסיבית (\textenglish{autoregressive}) של המודל
בזמן יצירה.


\section{קידוד מיקומי}

מאחר שמנגנון הקשב מחשב את ציוני הדמיון בין כל זוג אסימונים
ללא תלות בסדרם ברצף, הוא אינו מסוגל מטבעו לייצג את מיקומם
של האסימונים.
כדי לאפשר למודל לנצל מידע עמדתי, מוסיפים לייצוגי האסימונים
וקטורי קידוד מיקומי (\textenglish{positional encoding}) לפני
כניסתם למקודד ולמפענח.
גישה זו נחוצה הן לתמורות (permutations) הנדרשות בתחביר, שבהן
סדר המילים קובע את המשמעות, והן ליצירת טקסט קוהרנטי באורך
שרירותי.

\subsection*{קידוד סינוסואידלי}

וסמאני ועמיתיו \cite{vaswani2017attention} הציגו סכמת קידוד
מיקומי קבועה (\textenglish{fixed}) המבוססת על פונקציות
סינוס וקוסינוס בתדרים שונים:

\begin{equation}
  PE_{(pos, 2i)} = \sin\!\left(\frac{pos}{10000^{2i/d_{\text{model}}}}\right)
\end{equation}

\begin{equation}
  PE_{(pos, 2i+1)} = \cos\!\left(\frac{pos}{10000^{2i/d_{\text{model}}}}\right)
\end{equation}

כאשר \(pos\) הוא מיקום האסימון ברצף, \(i\) הוא אינדקס
הממד, ו-\(d_{\text{model}}\) הוא ממד הייצוג.
בחירת הפונקציות הסינוסואידליות אינה מקרית: לכל הזזת מיקום
\(\Delta pos\) קיימת הטרנספורמציה הלינארית המתאימה,
מה שמאפשר למודל להסיק מיקומים יחסיים ביעילות.
בנוסף, סכמה זו מרחיבה באופן טבעי לאורכי רצף שלא נצפו
בזמן האימון, שכן הפונקציות הסינוסואידליות מוגדרות עבור
כל ערך \(pos\) שהוא.

\subsection*{קידוד מיקומי נלמד}

בגישה חלופית, ניתן ללמוד את וקטורי הקידוד המיקומי
כפרמטרים של המודל.
מודל \textenglish{BERT} \cite{devlin2019bert} ומודל
\textenglish{GPT-2} \cite{radford2019language} אימצו גישה זו,
ולומדים מטריצת הטמעה (\textenglish{embedding matrix}) לעמדות
שממימדה \(L_{\max} \times d_{\text{model}}\), כאשר
\(L_{\max}\) הוא אורך הרצף המקסימלי שנקבע מראש.
יתרונה של גישה זו הוא גמישותה: המודל עשוי ללמוד ייצוגים
עמדתיים המותאמים לחלוטין לנתוני האימון.
עם זאת, חסרונה הוא שהמודל אינו יכול לעבד רצפים ארוכים
יותר מ-\(L_{\max}\) ואינו מכליל מחוץ לתחום האורכים שנראו
בזמן האימון.

\subsection*{קידוד מיקומי יחסי}

מחקר מאוחר של שוו ועמיתיו \cite{shaw2018relative} הציג
קידוד מיקומי יחסי (\textenglish{relative positional encoding}),
שבו המיקום המשולב בחישוב הקשב אינו המיקום האבסולוטי של
כל אסימון, אלא המרחק היחסי בין כל שני אסימונים.
גישה זו מקנה תכונת אינווריאנטיות להזזה (\textenglish{translation invariance}),
ומשפרת את הכללת המודל לאורכי רצף שונים.
בהמשך, ארכיטקטורת \textenglish{Transformer-XL} \cite{dai2019transformerxl}
הרחיבה רעיון זה ושילבה מנגנון זיכרון המאפשר עיבוד
רצפים ארוכים יותר מאלו שניתן לעבד בהעברה אחת.

כמכלול, שלושת הרכיבים שנבחנו בפרק זה — מנגנון הקשב
המסקאלרי, הקשב הרב-ראשי עם מבנה הקידוד-פענוח,
והקידוד המיקומי — מהווים את שלד הארכיטקטורה של
ה-\textenglish{Transformer} ואת הבסיס לכלל מודלי השפה
הגדולים שפותחו בעקבותיו, לרבות \textenglish{GPT-3}
\cite{brown2020language}, \textenglish{LLaMA} \cite{touvron2023llama},
ומודלים עבריים ייעודיים כגון \textenglish{AlephBERT}
\cite{seker2022alephbert} ו-\textenglish{HeBERT}
\cite{chriqui2022hebert}.

